\chapter{Future Work}\label{ch:future}

\section{Robust observation models and safe completion}
Three refinements follow from observed limitations of the model whose
results are supported on the held-out test sets. Each was evaluated on the
development set, and each is implemented as an optional capability that
leaves the reported model and its published tables unchanged. None of them
is part of the reported protocol, and adopting any of them would require a
further evaluation of the kind described in \S\ref{sec:hygiene}, fixed in
advance and carried out on fresh data. The three are a heavy-tailed
observation model for timing, a predictive floor as its cheaper alternative,
and a safe mode for the completion task.

A heavy-tailed (Student-$t$) observation model for timing works at both
levels at which it can be placed. As a deployment-time predictive,
variance-matched with $\nu = 5$, it removes the NLL tail demonstrated in
\S\ref{sec:guard-gp} completely and without any change to the fit: the
documented blow-up falls from $105.2$ to $-1.7$ with coverage intact. As an
inference-level noise model, a Gaussian scale mixture fit by leave-one-out
EM reweighting, it additionally brings timing coverage to nominal and
improves articulation recovery by $-0.018$, with an interval excluding zero.
The second effect arises because an observed timing outlier corrupts the
coupled channels through the coregionalization, and only a robust fit undoes
that corruption. Fits under the $t$ noise model must be scored with the
matched $t$ predictive, since a Gaussian score of a $t$ fit is incoherent.

Deployment-time predictive floors, the seemingly cheap engineering
alternative, were found to be ineffective. The failure is a tail-shape
problem rather than a scale problem. A floor raises the predictive scale of
every note, but the outliers that dominate the tail already carry an
appropriate predictive scale, so a floor cannot reach them. The negative
result is useful in that it rules out the simplest remedy and locates the
problem in the shape of the likelihood.

The completion task of \S\ref{sec:boundary} can be made safe against
extrapolation by a combination of two tests, each of which is only half of
the answer on its own. A Mahalanobis coverage test, workable only in the
observed excerpt's low-rank frame, is silent under interpolation, as
required, but blind to catastrophic notes that lie inside the excerpt's
feature coverage. A $3\sigma$ disagreement test against the cross-piece head
supplies the other half. Together they tame every blow-up, the worst
evaluation falling from $480{,}266$ to $0.7$, while keeping the per-piece
model's advantage wherever adaptation helps. Blow-ups also occur under
interpolation masks, so the boundary concerns the determination of the
feature weights rather than extrapolation alone.

\section{Strings and winds: open questions}
The strings-and-winds result of Chapter~\ref{ch:phase2} leaves four open
items: the corpus, the timing channel, the tonal metric, and additional
channels. The first concerns corpora beyond URMP. In URMP the note onsets
are annotated, and the timing channel's noise row is the predictive variance
of a local tempo-line fit to those annotations. On any corpus without
note-level onset annotations the general audio-alignment problem returns,
and the timing channel's noise row must then carry an aligner's error.

The timing hypothesis of the first strings-and-winds evaluation, the
calibration of $\tau$, was not supported, and the diagnosis named its own
follow-up: a correlated noise row for $\tau$ that represents the structure
of alignment error along score time. That follow-up was first evaluated on
the development set. An AR(1) block over note order, with its one parameter
chosen by the evidence, improves timing NLL by $-0.111$ with an interval
excluding zero, moves coverage toward nominal, and reduces timing RMSE by
${\sim}20\%$ ($\approx 99 \to 80$\,ms). The correlation was detected on
97--100\% of the piece-seed combinations across three independent runs,
with fresh mask seeds and with a joint refit. The revised model was then
supported on the second, disclosed evaluation of the held-out test set
(\S\ref{sec:phase2-reuse}): the paired differences are $-0.121$ in timing
NLL and $-0.011$ s in timing RMSE, both with intervals excluding zero, and
coverage is at nominal. What remains is adoption. Making the correlated row
the default noise model of the timing channel would change the published
tables and therefore calls for its own evaluation fixed in advance, and on a
fresh corpus the hypothesis is untested.

The circle-of-fifths replacement geometry \eqref{eq:tonaladj}, which hurt
piano expression, helps intonation on both axes in the development set
(\S\ref{sec:phase2-delta}). It was evaluated a second time on the held-out
test set with calibration as the pre-specified primary measure
(\S\ref{sec:phase2-reuse}). The recovery gain held, $-0.312$ cents with an
interval excluding zero, and the calibration gain did not, $-0.023$ with an
interval reaching zero. The metric therefore stays exploratory. What remains
is a fresh corpus and a design powered for recovery, which a power analysis
made before that evaluation places above ${\sim}20$ pieces. Nothing from
this route is adopted.

The within-note drift and decay slopes measured in \S\ref{sec:phase2-delta}
remain candidates for additional channels. They would follow the template of
the vibrato onset delay $d^{\mathrm{vib}}$, which entered the bundle as an
output with its own estimator variance and with no prediction attached to
it. Neither slope is part of the reported bundle, and adding either would
extend the output rather than change the model. Whether they are worth
adding depends on whether a transcription use-case calls for them.

\section{The waveform: the joint loop}
\label{sec:future-phase3}Chapter~\ref{ch:phase3} stops at chains of a few
notes. Within those chains the score prior couples the per-note constants
(\S\ref{sec:phase3-priorvalue}), where it transfers only the context's
common offset, and the curves themselves still carry independent priors.
What remains is the \emph{joint} loop, in which all notes of a performance
are inferred at once under the graph prior, by Laplace or variational
inference over $z$. Three further components belong to that loop: a
transient model at the note joints, which the chain study indicates as the
next step, the estimand bridge between the harmonic model's constant and the
annotation chain's functional, and the amplitude envelope $A_i(t)$ brought
under the prior. The evaluation design changes with it. With no per-note
observations there is no visible/hidden split, and the natural protocol is
held-out pieces with posterior-predictive resynthesis.

\section{Incorporating music theory}\label{sec:musictheory}
Music theory can inform the model in two places, the graph or kernel and
the prior mean, which is equivalently the feature kernels. The two routes
are not exclusive, each is testable by the same held-out
recovery-and-calibration criterion, and both have been evaluated. The graph
route gave one positive and one negative result, both instructive, and the
feature route gave a negative result that an embedding probe explains. A
single criterion, stated at the end of this section, decides whether any
ingredient from either route is kept.

\subsection{In the graph and the kernel}
Two constructions in this route are implemented (\S\ref{sec:graphvariants})
and evaluated (\S\ref{sec:kernels}), and they gave opposite results.
\emph{Replacing} the semitone ruler with a circle-of-fifths pitch metric
\eqref{eq:tonaladj} makes both recovery and calibration worse, with paired
intervals excluding zero: expressive behaviour travels with register
proximity, not tonal proximity. \emph{Adding} same-onset (chord) and
stepwise voice-leading edges \eqref{eq:harmadj} on top of the plain graph is
the one intervention that improves both axes with intervals excluding zero.
The general lesson, with the qualifier that the strings-and-winds setting
supplies, is that music theory earns its place in the kernel as extra
relations rather than as a substitute geometry, for channels whose target
is not pitch. On intonation, where the target is pitch itself, the same
substitute geometry helps (Chapter~\ref{ch:phase2}), so the rule is
channel-relative.

What remains open in this route is the part that the score alone cannot
supply. Cadential and functional edges, such as tonic and dominant
relations, require a harmonic analysis, since the implemented chord edges
encode onset simultaneity and not chord identity
(\S\ref{sec:graphvariants}). Metrical-hierarchy and phrase weights could
form further edge families. Their feature-level counterparts were evaluated
in the feature-kernel route below, but as graph structure they are untested.
A Tonnetz or neo-Riemannian lattice metric and the Tymoczko voice-leading
distance are alternatives to \eqref{eq:fifths} that have not been tried.

\subsection{In the feature kernels}
Fourteen music-theoretic score columns, covering local key, metrical weight,
dissonance, phrase salience, motif repetition and tessitura, change nothing
on any channel at any masking level, with no paired interval excluding zero.
A linear and random-features probe of the embeddings rules out the obvious
explanation, that the columns duplicate what the music model already
encodes. The embeddings encode voicing, rhythm and register but \emph{not}
tonality: scale degree, mode and local key sit at chance. Taken with the
tonal-metric result of the graph route, the consistent reading is that tonal
structure carries no measurable marginal signal for these three expressive
channels at this scale, while the rhythmic and voicing structure that the
columns add is already in the kernel. The full study is reported in
Appendix~\ref{app:theoryfeat}. The columns remain available as an
interpretable probe and are not part of the reported model.

A single criterion decides whether an ingredient from either route is kept.
A music-theory ingredient is retained only if it measurably improves
recovery or calibration over the plain score graph on held-out data.
Otherwise it is reported as motivation rather than as machinery. Both routes
are directly testable within the existing protocol, so the criterion can be
applied to every candidate listed above without a new evaluation design.

\section{Unknown score support}\label{sec:future-amt}
A discrete latent layer over note identities and onsets would move the model
toward automatic music transcription proper. The score support $S$, assumed
known throughout this thesis, would become a latent variable with its own
prior, and the expressive variables would be inferred jointly with it. Every
component of the present model is conditioned on $S$, so the model would be
retained as the conditional inference inside a larger inference over the
score. It is the assumption relaxed last, once the structured performance
model is established.

% =====================================================================
